% Source: 38_Chapter_28_Another_Man_At_Manresa_The_Catholic_Ref.txt
\bookchapter[“Another Man” at Manresa]{CHAPTER 28}{“ANOTHER MAN” AT MANRESA}{The Catholic Reformation}

\leadin{I}{N 1521, THE YEAR} that Martin Luther stood before the Emperor Charles V at the Diet of
Worms, a young Spanish nobleman was fighting on the emperor’s borderlands against the invading
French at Pamplona. A cannonball shattered one of his legs. During a long, painful convalescence he
turned, out of boredom, to two popular inspirational works, one on the lives of the saints and the
other a life of Christ. With these his long process of conversion began.

Months later, at the Benedictine abbey of Montserrat, he exchanged his gentleman’s clothes for a
rough pilgrim’s garb and dedicated his sword and dagger to the shrine’s Black Virgin. For nearly a
year, in a little town called Manresa, thirty miles north of Barcelona, he gave himself to an orgy
of austerity: begging door to door, wearing a barbed girdle, fasting for days on end. For months he
endured the terrible depressions of the mystic’s dark night of the soul, even contemplating suicide
at one point. But what followed was the mystic’s singular reward, an immense breakthrough to
spiritual enlightenment. In a wave of ecstatic illumination one day at the River Cardoner, the
wounded nobleman, Ignatius Loyola, became, in his own words, “another man.”

Loyola (1491–1556) transformed his rebirth at Manresa to a plan for spiritual discipline, a military
manual for storm troopers at the service of the pope. The result was the Jesuits, the Society of
Jesus, the greatest single force in Catholicism’s campaign to recapture the spiritual domains seized
by Protestantism.

How did the Church of Rome respond to the Protestant challenge? It didn’t, not immediately. But when
it finally realized the seriousness of the revolt, it called upon its spiritual warriors, it
convened a new militant council, and it reformed the machinery of the papal office. Faced by the
rebellion of almost half of Europe, Catholicism rolled back the tide of Protestantism until by the
end of the sixteenth century Protestantism was limited roughly to the northern third of Europe, as
it is today.

Some historians have interpreted the Catholic Reformation as a counterattack against Protestantism;
others have described it as a genuine revival of Catholic piety with few thoughts of Protestantism.
The truth is the movement was both a Counter Reformation, as Protestants insist, and a Catholic
Reformation, as Catholics argue. Its roots run back to forces before Luther’s time, but the form it
took was largely determined by the Protestant attack.

\begin{sidebar}{Facets of the Counter Reformation}
\begin{numlist}
\item[1.] Ignatius became an important voice for Catholic spirituality, as well as organizing the
  influential Jesuit order.
\item[2.] The Counter Reformation also contributed to the emergence of other important voices in
  spirituality, including St. Teresa of Avila.
\item[3.] Catholics answered the academic and scholarly work of the Protestants with work in
  language (especially in Spain) and training for priests.
\item[4.] The council of Trent offered an articulation of doctrines in reaction to the Protestant
  teaching.
\item[5.] The Catholic Church initiated a global missionary expansion.
\item[6.] The Inquisition, an adaptation of an ancient Roman legal format, proved an effective yet
  brutal strategy to discover and stop heresy.
\end{numlist}
\end{sidebar}

\section{RETURN TO THE SPIRITUAL}

Strange as it may seem the mystical experience was a large part of Catholicism’s recovery. The
sixteenth century produced a remarkable variety of Catholic saints: the English lawyer and
statesman, Thomas More; the cheerful and imaginative missionary to the Calvinists, Francis of Sales;
the somber reforming archbishop of Milan, Charles Borromeo; the rapturous Spanish mystic, Teresa of
Avila; and, most influential of all, the Spanish soldier of Christ, Ignatius Loyola.

Even before Luther posted his theses on the church door, a distinguished and aristocratic group at
Rome had formed a pious brotherhood called the Oratory of Divine Love. Their guiding belief was that
the reformation of the church and society begins within the individual soul.

The Oratory was never large in number, perhaps fifty, yet it had enormous influence. It stimulated
reform in the older monastic orders and contributed leaders to the Church of Rome as it laid plans
for a general council to deal with internal reform and the Protestant heresy. Among the members of
the Oratory who later emerged as significant figures were Jacopo Sadoleto, who debated with Calvin;
Reginald Pole, who tried under Bloody Mary to turn England back to Rome; and Gian Pietro Caraffa,
who became Pope Paul IV.

Throughout the 1520s and 1530s, however, the Church of Rome took no significant steps toward reform.
The question is why. Why was she so slow to respond to the Protestant challenge?

One simple answer is politics. The Emperor Charles V and the popes fought a running battle over the
calling of a general council that stretched over two decades. Luther had called for a council of the
church as early as 1518. The idea gained the support of the German princes and the emperor, but the
popes had fears of such an assembly. They remembered too well the councils at Constance and Basel.
They also knew that many in Germany had in mind a council without the pope.

Equally important, the popes in the 1520s and 1530s were preoccupied with secular and political
affairs. Clement VII (1523–1534) is a prominent example. He regarded concern for the Papal States in
Italy as a supreme law, and his passion for papal political fortunes drew him into an alliance with
France against Charles V, leader of the Hapsburg interests in Italy. The pope’s treachery and
disloyalty enraged Charles, and he threatened Clement with a trial before a general council unless
he broke his alliance with Francis I, King of France.

To show the pope he meant business Charles ordered his troops to march on Rome. As it turned out he
got more than he had planned. The commanders of his troops were killed. As a result, the rough and
undisciplined Spanish and German mercenaries were leaderless when they stormed Rome on May 6, 1527.
Their pillaging, plundering, and murdering in the Eternal City lasted for weeks. The pope took
refuge in the Castle of St. Angelo, but finally had to surrender and endure half a year of harsh
imprisonment. Many considered this sack of Rome a terrible visitation by God, a clear call to
repentance and change in the worldly papacy.

No serious reform came until Pope Paul III (1534–1549) ascended the papal throne. Paul appeared to
be a most unlikely candidate for spiritual leadership. He had three illegitimate sons and a
daughter, four striking reminders of his pursuit of pleasure. The sack of Rome, however, seems to
have sobered him. He realized that the time had come for reform to begin in the House of God. He
started where he felt a change of heart was most urgently needed, in the College of Cardinals. He
appointed to the college a number of champions of reform. Among them were leaders in the Oratory of
Divine Love: Sadoleto, Pole, and Caraffa.

Paul then appointed nine of the new cardinals to a reform commission. The head of the commission was
another former member of the Oratory, Gasparo Contarini. A peacemaker by temperament, Contarini
stood for reconciliation with the Protestants and advocated a return to the faith of the apostles.

After a wide-ranging study of conditions in the Church of Rome, the commission issued in 1537 a
formal report, Advice~.~.~. Concerning the Reform of the Church. Disorder in the Church, the report
said, could be traced directly to the need for reform. The papal office was too secular. Both popes
and cardinals needed to give more attention to spiritual matters and stop flirting with the world.
Bribery in high places, abuses of indulgences, evasions of church law, prostitution in Rome: these
and other offenses must cease.

\section{CALL FOR A GENERAL COUNCIL}

Pope Paul took action on a few of these items, but his most significant response was a call for a
general council of the Church. After intense negotiations he agreed with the emperor on a location
for the assembly, a town in northern Italy under imperial control called Trent.

Even then, however, no council assembled for years because Francis I did everything in his power to
prevent such a gathering. In his passion for leadership of Europe, he feared a council would only
strengthen Charles’s hand. Francis even stooped to incite the Turks against the emperor. Two wars
between Francis and Charles delayed the opening of a council until 1545, almost three decades after
Luther’s theses appeared.

By 1545 Rome was under the spell of a new austerity. Reform was on the rise. The immoralities of
Paul’s younger days were no longer acceptable behavior. The pope’s new rigor was apparent in the
institution of the Roman Inquisition and in the Index of prohibited books: works that any Catholic
risked damnation by reading. All the books of the Reformers were listed, as well as Protestant
Bibles. For a long time merely to possess one of these banned books in Spain was punishable by
death. The Index was kept up to date until 1959 and was finally abolished by Pope Paul VI.

\section{IGNATIUS LOYOLA}

It was this city, a new militant Rome, and this pope, Paul III, who approved Ignatius Loyola’s new
Society of Jesus. These daring soldiers of Christ promised the pope that they would go wherever he
might send them “whether to the Turks or to the New World or to the Lutherans or to others, be they
infidel or faithful.” They said it and they meant it, for the Jesuits were a fitting likeness of
their founding general.

As a boy Ignatius had left the gloomy castle of Loyola near the Pyrenees to enter the court of his
father’s noble friend. He had grown into little more than an engaging playboy, spending his days in
military games or reading popular chivalrous romances, his nights pursuing less noble adventures
with local girls.

But all that was before Ignatius met God at Manresa. Martin Luther emerged from his spiritual
struggle convinced that the human will is enslaved, that man cannot save himself. God, and God
alone, must deliver him. Loyola came out of his struggle believing that both God and Satan are
external to man, and man has the power to choose between them. By the disciplined use of his
imagination man can strengthen his will so as to choose God and his ways.

One of Loyola’s spiritual exercises, for example, aimed to make the horrors of hell real:

\begin{blockquote}
Hear in imagination the shrieks and groans and blasphemous shouts against Christ our Lord and all
the saints. Smell in imagination the fumes of sulphur and the stench of filth and corruption. Taste
in imagination all the bitterness of tears and melancholy and growing conscience. Feel in
imagination the heat of the flames that play on and burn the souls.
\end{blockquote}

\noindent The same technique, of course, could be used to represent the beauties of the Nativity or
the glories of heaven. By proper discipline the imagination could strengthen the will and teach it
to cooperate with God’s grace.

For Ignatius personally surrender to the will of God meant more education. He entered a school in
Barcelona to sit with boys less than half his age to study Latin, then threw himself into a dizzying
year of courses at the University of Alcala. Out of it came his conviction that learning must be
organized to be useful. The idea eventually grew into the Jesuits’ famed plan of studies, which
measured out heavy but manageable doses of classics, humanities, and sciences.

Ignatius became such a fervent evangelist that the Inquisition imprisoned and examined him more than
once about his life, teaching, and theology. Perturbed, he left for Paris where he spent seven years
at the university, became “Master Ignatius,” and gathered around him the first of his permanent
companions: Peter Faber, Diego Laynez, Alfonso Salmeron, Simon Rodriguez, Nicholas Bobadilla, and,
above all, the young Spanish nobleman, Francis Xavier.

Ignatius shared with them his remarkable program for sainthood, his Spiritual Exercises. A
distillation of his own religious experiences during and following his conversion, the Exercises
prescribe four “weeks” of meditations, beginning with sin, death, judgment, and hell, and then
moving on to Christ’s life, death, and resurrection.

Ignatius intended a path to spiritual perfection: rigorous examination of conscience, penance, and a
resolute amnesia about guilt once the spiritual pilgrim confronted God’s forgiveness. The Exercises
became the basis of every Jesuit’s spirituality. Later popes also prescribed them for candidates for
ordination, and Catholic retreats applied them to lay groups.

In his Portrait of the Artist as a Young Man, James Joyce describes his young hero, Stephen Dedalus,
reduced to horror by a sermon on hell: “A wave of fire swept through his body~.~.~. flames burst
forth from his skull.” After he had gone to confession, however, “the past was past.” That is just
the way Ignatius planned it. “Send no one away dejected,” Loyola wrote. “God asks nothing
impossible.” Thus his followers became the great apostles of the possible.

\section{THE SOCIETY OF JESUS}

In 1540 Pope Paul III approved the little Society of Jesus as a new religious order. In Ignatius’s
metaphor, they were to be chivalrous soldiers of Jesus, mobile, versatile, ready to go anywhere and
perform any task the pope assigned. As a recognized order, they added to their earlier vows of
poverty and chastity the traditional vow of obedience to their superiors and a fourth vow expressing
their special loyalty to the pope. They gave command to a superior general elected for life. Their
choice for the first general was Ignatius.

The aim of the order was simple: to restore the Roman Catholic Church to the position of spiritual
power and worldly influence it had held three centuries before under Innocent III. Everything was
subordinated to the Church of Rome because Ignatius believed firmly that the living Christ resided
in the institutional church exclusively.

Perhaps the most fascinating feature of the Jesuits was their perilous attempt to live energetically
in the world without being of it. Loyola wanted them to be all things to all men, and they nearly
succeeded.

Their efforts have rallied defenders who assign them the highest posts in heaven and critics who
consign them to the lower regions of hell. Writing to Thomas Jefferson in 1816, John Adams said, “If
ever any congregation of men could merit eternal perdition on earth and in hell, it is the company
of Loyola.” Whatever their destination, the Jesuits were from their beginning unique.

That first generation under Loyola’s zealous leadership rode full gallop into their new assignments:
convert the heathen, reconvert Protestant Europe. Francis Xavier leaped from India to Southeast Asia
to Japan, a country that had never before heard the Christian message. More than any others, the
Society of Jesus stemmed, and sometimes reversed, the tide of Protestantism in France, the Low
Countries, and Central Europe. When Ignatius died in 1556, his order was nearly one thousand strong
and had dispatched its apostles to four continents.

No mission in that first generation proved more decisive than the band of Jesuits assigned to the
Council of Trent. Only thirty-one council fathers led by three papal legates were present for the
opening ceremonies of the council. None of them could have guessed that their modest beginning would
lead to the most important council between Nicaea (325) and Vatican II (1962–1965). Under the
influence of the Jesuits, Trent developed into a powerful weapon of the Counter Reformation. Two
suave, intelligent, and highly influential members of the society, Diego Laynez and Alfonso
Salmeron, guided the agenda more and more toward “the correct churchly attitude” of the followers of
Loyola.

The council fathers met in three main sessions: 1545–1547, 1551–1552, and 1562–1563. Throughout the
sessions the Italians were strongly represented. Other areas, notably France, were significantly
underrepresented. Compared to other councils, Trent was never well attended. During the second
series of sessions a number of Protestants were present, but nothing came of it. From start to
finish the council reflected the new militant stance of Rome.

\section{THE SHAPE OF MODERN CATHOLICISM}

Everything the Protestant Reformation stood for was vigorously—one could almost say
violently—rejected at Trent. The Protestant Reformers emphasized justification by faith alone. The
council rejected a potential compromise holding that justification is the result of two causes: 1) a
righteousness that is external to the disciple but added to our account in the eyes of God, and 2) a
righteousness that is internal, infused into the disciple. The Protestant’s external righteousness
was an event whereby God counted the person right in his eyes. For Catholics internal righteousness
was the process whereby God imparted righteousness within the disciple during the course of their
life. Trent emphatically linked justification to the process of becoming righteous, leaving the
notion of being declared righteous to the Protestant proclaimers.

Luther, Calvin, and Grebel stressed salvation by grace alone; the council emphasized grace and human
cooperation with God to avoid, in Loyola’s terms, “the poison that destroys freedom.” “Pray as
though everything depended on God alone;” Ignatius advised, “but act as though it depended on you
alone whether you will be saved.”

The Protestants taught the religious authority of Scripture alone. The council insisted on the
supreme teaching office of the Roman Church, popes and bishops, as the essential interpreters of the
Bible.

Thus the Council of Trent guaranteed that modern Roman Catholicism would be governed by the
collaboration between God and man. The pope remained, the seven sacraments remained, the sacrifice
of the mass remained. Saints, confessions, indulgences remained. The council’s work essentially
narrowed the range of options previously available to medieval Catholics; the anti-Protestant
reactionary tone was unmistakable. After four centuries, we can look back to the Reformation age and
see that the religious unity of western Christendom was permanently shattered. Men and women in
Loyola’s lifetime did not see that truth. The fact dawned upon Europe slowly.

At first Luther’s followers thought him so obviously right that the Catholic Church would inevitably
adopt his ideas. Others thought him so obviously wrong that sooner or later he would be burned as a
heretic and his movement wither away. After all, the past was strewn with the corpses of heretical
bodies. The point is both parties, Catholic and Protestant, thought that they represented the true,
catholic church of Christ and their enemies a false version of it. That was the spirit of Trent as
well.

As time went on, however, the thinking of ordinary men reached, almost imperceptibly, a second
stage. The uneasy, half-conscious conviction grew that the conflict had reached a stalemate.
Catholicism could not crush the new heresy, and Protestantism could not overthrow Rome. In this
second stage men had no real emotional or intellectual acceptance of a stalemate, only a bitter
admission of the fact.

The vast majority of people were still convinced that religious truth was identifiable. Truth stood
on one side, error on the other. Error meant not only individual damnation, but infection of others
and destruction of society. Resistance of these dreadful evils took the form of inquisitions, civil
wars, and persecutions.

In this ideological warfare, Calvinists and Jesuits led the charge. Each embodied a militant
organization calling for loyalty above and beyond national and political ties. At this stage almost
no one could imagine that truth could lie on both sides of the battlefield or that both parties
could peacefully coexist in the same state or even on the same continent.

The third stage of religious toleration, based upon the full acceptance of religious diversity
within a nation, appeared before 1600 only in hints in the attitudes of mystics, Christian humanists
like Erasmus, radical Protestants like the Anabaptists, and practical politicians like Queen
Elizabeth of England.

Perhaps more powerful still were the political and economic forces committed to the emergence of the
nation-state. Europe now had small political areas (the size of large counties or small states in
the U.S.) under the umbrella of an emperor and the Catholic Church. These smaller realms under lords
were merging into what we recognize today as nations, such as Germany. A state with geographic
boundaries and more people would require both Protestant and Catholic citizens to coexist and
cooperate.

Some voices, anticipating later Enlightenment thinking, called for a generic concept of God shared
by virtually all believers in the West to become a new public standard: God is eternal, one, etc.
Distinctive beliefs and social practices would be tolerated if they could be transformed to become a
matter of private devotion. Feeble as it was, this whisper was the voice of the future. No company
of Christians resisted it more doggedly than the Jesuits.

\readinghead{Suggestions for Further Reading}
\begin{readinglist}
\entry *Bireley, Robert. The Refashioning of Catholicism, 1450–1700. Washington, D.C.: Catholic
  University of America Press, 1999.
\entry Brodrick, James. The Origin of the Jesuits. London: Longmans, Green and Co., 1949.
\entry Daniel-Rops, Henry. The Catholic Reformation. New York: E. P. Dutton, 1962.
\entry Franzen, August. A History of the Church: Revised and Edited by John P. Dolan, New York:
  Herder and Herder, 1969.
\entry Noll, Mark A. Turning Points: Decisive Moments in the History of Christianity. Grand Rapids:
  Baker Books, 2000.
\entry *O’Malley, John W. The First Jesuits. Cambridge, MA: Harvard University Press, 1993.
\entry Thompson, Francis. Saint Ignatius Loyola. Westminster, MD: Newman Press, 1950.
\entry Wakefield, James L. Sacred Listening: Discovering the Spiritual Exercises of Ignatius Loyola.
  Grand Rapids: Baker Books, 2006.
\end{readinglist}
